% Generated draft: requires mathematical and source review.
\documentclass[11pt]{article}
\usepackage[margin=25mm]{geometry}
\usepackage{fontspec}
\setmainfont{Latin Modern Roman}
\usepackage{amsmath,amssymb,mathtools,unicode-math}
\setmathfont{Latin Modern Math}
\usepackage{enumitem,needspace,xurl,hyperref}
\setlength{\parindent}{0pt}
\setlength{\parskip}{5pt}
\setlength{\emergencystretch}{3em}
\newcommand{\N}{\mathbb{N}}
\newcommand{\Z}{\mathbb{Z}}
\newcommand{\Q}{\mathbb{Q}}
\newcommand{\R}{\mathbb{R}}
\newcommand{\C}{\mathbb{C}}
\newcommand{\F}{\mathbb{F}}
\newcommand{\A}{\mathbb{A}}
\newcommand{\PP}{\mathbb{P}}
\newcommand{\E}{\mathbb{E}}
\newcommand{\eps}{\varepsilon}
\newcommand{\dd}{\,\mathrm{d}}
\begin{document}
\section*{UnsolvedMath ID: 10600004}
\textbf{Source problem code:} catalogue entry
\subsection{Definitions and mathematical statement}
A planar doodle is represented here by a generic immersed collection of closed curves in the plane, with transverse double points recorded as flat crossings, modulo planar isotopy and flat Reidemeister moves of types I and II. Type III is forbidden; equivalently, a deformation is not allowed to pass through a triple point. A flat virtual knot diagram instead has flat crossings and distinguished virtual crossings, with the flat virtual equivalence moves of the intended theory. In particular, allowing the flat type III move changes the equivalence relation from that of doodles.

The problem asks for a comparison of flat virtual knot theories with doodle theories, including the version of doodles with virtual as well as flat crossings. The latter is described as giving doodles on surfaces of higher genus; its complete move system and the permitted surface equivalences or stabilizations are not specified in the extract. A precise comparison must fix these categories and determine how their objects, equivalence relations and invariants correspond; the source supplies no single proposed classification or comparison theorem.

The accompanying discussion reports that Khovanov associated a group to a doodle, that commutator identities can also be associated to doodles, and that doodle cobordism is defined by an immersed surface without triple points. It says the resulting cobordism classes represent elements of a free abelian group. The construction of that group, its generators or rank, the group associated by Khovanov, and the cobordism orientation and boundary conditions are not given. These statements are preserved as source context for the requested comparison, without treating them as definitions of an unspecified invariant.
\subsection{Short English statement}
Compare flat virtual knot theories with planar and virtual doodle theories, retaining the essential prohibition of the flat type III move and of triple points in doodle cobordisms.
\subsection{Sources}
\begin{itemize}
\item \emph{Unsolved Problems in Virtual Knot Theory and Combinatorial Knot Theory}, revised survey, dedicated problem-list item 4, ``Virtuals and the Theory of Doodles.'' \url{https://arxiv.org/abs/1409.2823}. Recorded access 2026-07-29; extraction source-tex. Author names and a specific revision identifier are not supplied in this record.
\item The statement attributes a doodle group to Khovanov, but its reference is missing from the extracted text; the final higher-genus ``see'' reference is also missing.
\item Local UnsolvedMath registry, \texttt{data/unsolvedmath/problems.json}, permanent ID \(10600004\), exact source code \texttt{catalogue entry}. Supplied status partially\_solved and research classification PARTIAL-PROGRESS; the summary mentions active virtual-doodle and special-setting invariant results without precise references.
\end{itemize}
\end{document}
